\chapter{Antes de instalar}\label{cap:02-antes-de-instalar}

Este capítulo prepara todo lo necesario: una copia de tus archivos, la
descarga de Lázaro y la memoria USB desde la que arrancará el equipo.
Calcula una hora, casi toda esperando a que termine la descarga.

Vas a necesitar:
\begin{itemize}
  \item el equipo donde quieres instalar Lázaro, enchufado a la corriente;
  \item otra computadora con internet para descargar (puede ser la misma,
        si todavía funciona con Windows);
  \item una memoria USB de 4 GB o más (8 GB para la edición Completa) que
        no tenga nada importante: \textbf{se va a borrar}.
\end{itemize}

\section{Guarda una copia de tus archivos}\label{sec:copia-antes}

\begin{cuidado}
Al instalar Lázaro se borra el disco que elijas, con todo lo que tiene. No
hay forma de recuperar lo borrado. Haz la copia \textbf{antes} de seguir.
\end{cuidado}

Lo que suele importar está en estas carpetas de Windows: \emph{Documentos},
\emph{Imágenes}, \emph{Música}, \emph{Videos}, \emph{Descargas} y
\emph{Escritorio}. Cópialas a un pendrive o a un disco externo (no al
mismo pendrive que usarás para Lázaro). Revisa también:
\begin{itemize}
  \item \textbf{Las contraseñas y marcadores del navegador.} Si usas Chrome
        o Edge con una cuenta, ya están guardados en esa cuenta. Si no,
        expórtalos: en Chrome, \menu{Marcadores > Administrador de marcadores
        > \ldots > Exportar marcadores}.
  \item \textbf{El correo.} Si lo lees en una página web (Gmail, Outlook.com,
        Yahoo) no tienes que hacer nada: sigue en internet.
  \item \textbf{Las fotos del celular} que hayas pasado al equipo.
  \item \textbf{Claves de wifi}: anota la de tu casa; la vas a necesitar.
\end{itemize}

Para hacer una copia de seguridad bien hecha y conservarla en el tiempo,
el capítulo~\ref{cap:11-archivos} da más detalles. El Instituto Nacional de
Ciberseguridad de España tiene una buena guía para no expertos
\cite{incibe-copias}.

\section{Descarga Lázaro}\label{sec:descargar}

Primero elige la edición (sección~\ref{sec:ediciones}): \textbf{Ligera} si
tienes buena conexión, \textbf{Completa} si el equipo no va a tener
internet. Luego descarga \emph{dos} archivos de la misma carpeta:

\begin{itemize}
  \item el archivo que termina en \texttt{.iso}: es el \termino{archivo ISO}, un
        \enquote{disco} con todo el sistema dentro. Pesa entre 2 y 4,6 GB;
  \item el archivo del mismo nombre que termina en \texttt{.sha256}: es muy
        pequeño y sirve para comprobar que la descarga llegó bien.
\end{itemize}

Los dos están en:
\begin{itemize}
  \item \textbf{Ligera:} \url{https://github.com/Entropio1984/lazaro-linux/releases}
        o \url{https://sourceforge.net/projects/lazaro-linux/files/Ligera/}
  \item \textbf{Completa:} \url{https://sourceforge.net/projects/lazaro-linux/files/Completa/}
  \item \textbf{Versiones numeradas} (fijas, las más probadas):
        \url{https://sourceforge.net/projects/lazaro-linux/files/Versiones/}
\end{itemize}

No hace falta crear ninguna cuenta. Si la descarga se corta, vuelve a
empezarla: el navegador suele poder continuarla desde donde quedó.

\section{Comprueba que la descarga llegó entera}\label{sec:sha256}

Una descarga dañada es la causa más común de que el USB no arranque, y
tiene arreglo fácil si se detecta a tiempo. El archivo \texttt{.sha256}
trae una \termino{suma de comprobación}: un número largo que se calcula a
partir de cada byte de la ISO. Si cambia un solo byte, cambia el número
entero.

\subsection*{En Windows}
\begin{pasos}
  \item Abre la carpeta donde quedó la descarga (normalmente \emph{Descargas}).
  \item Haz clic derecho en un espacio vacío de la carpeta y elige
        \menu{Abrir en Terminal} (en Windows 10: mantén \keys{\shift}
        apretada mientras haces clic derecho y elige \menu{Abrir la ventana
        de PowerShell aquí}).
  \item Escribe esto y pulsa \keys{Enter}:
\begin{terminal}
Get-FileHash .\textbackslash lazaro-linux-*.iso
\end{terminal}
  \item Tarda un minuto. Aparece una línea con el número (la columna
        \emph{Hash}). Abre el archivo \texttt{.sha256} con el Bloc de notas y
        compara: deben ser iguales. Basta con mirar los primeros y los
        últimos 6 caracteres \cite{ms-gethash}.
\end{pasos}

\subsection*{En Linux (otra distribución, o Lázaro ya instalado)}
Abre una terminal en la carpeta de la descarga y escribe:
\begin{terminal}
sha256sum -c lazaro-linux-*.sha256
\end{terminal}
Debe decir \texttt{La suma coincide} o \texttt{OK}.

\begin{cuidado}
Si los números \textbf{no coinciden}, la descarga está dañada: bórrala y
descárgala de nuevo, mejor desde la otra página (GitHub o SourceForge). No
sigas con una ISO dañada: el equipo podría no arrancar o fallar a mitad de
la instalación.
\end{cuidado}

\section{Prepara la memoria USB}\label{sec:usb}

Grabar la ISO en el USB no es lo mismo que copiarla: hace falta un programa
que la convierta en una memoria \emph{arrancable}. Elige uno de estos tres,
todos gratuitos.

\begin{cuidado}
Esto borra todo lo que tenga la memoria USB.
\end{cuidado}

\subsection{balenaEtcher (Windows, Mac y Linux): el más sencillo}
\begin{pasos}
  \item Descárgalo de \url{https://etcher.balena.io} e instálalo
        \cite{etcher}.
  \item Ábrelo y pulsa \menu{Flash from file}. Elige la ISO de Lázaro.
  \item Pulsa \menu{Select target} y elige tu memoria USB. Fíjate bien en
        el tamaño para no elegir otro disco.
  \item Pulsa \menu{Flash!} y espera. Al final Etcher comprueba la copia
        solo.
\end{pasos}

\subsection{Rufus (Windows)}
\begin{pasos}
  \item Descárgalo de \url{https://rufus.ie/es/}: no necesita instalarse
        \cite{rufus}.
  \item En \emph{Dispositivo}, elige tu memoria USB.
  \item Pulsa \menu{Seleccionar} y elige la ISO de Lázaro.
  \item Pulsa \menu{Empezar}. Si pregunta cómo escribir la imagen, elige
        \textbf{Escribir en modo Imagen DD}. Ese modo copia la ISO tal cual
        y es el que funciona con Lázaro.
\end{pasos}

\subsection{Ventoy: varias ISO en un mismo USB}
Ventoy prepara el USB una sola vez; después basta con copiar los archivos
ISO dentro, como a cualquier pendrive, y al arrancar te deja elegir cuál
usar \cite{ventoy}. Es cómodo si quieres llevar Lázaro junto a otros
sistemas.

\begin{cuidado}
Con Ventoy, después de copiar la ISO usa siempre \menu{Quitar el
dispositivo de forma segura} antes de sacar el USB. Si lo sacas de golpe,
la copia puede quedar dañada y el equipo no arranca o se detiene con un
\enquote{error de entrada/salida}. Copia también el archivo \texttt{.sha256} al
USB: así puedes repetir la comprobación sobre la copia que quedó en él.
\end{cuidado}

\section{BIOS, UEFI y el arranque seguro}\label{sec:bios-uefi}

Cuando se enciende, todo equipo ejecuta primero un pequeño programa propio
que revisa el hardware y busca desde dónde arrancar. En los equipos de
antes de 2012, aproximadamente, se llama \termino{BIOS}. En los más nuevos
es \termino{UEFI}, su reemplazo moderno \cite{uefi-org}. Lázaro funciona con
los dos y no tienes que hacer nada especial.

Muchos equipos con UEFI traen activado el \termino{arranque seguro}
(\emph{Secure Boot}), que solo deja arrancar sistemas firmados con llaves
de Microsoft \cite{ms-secureboot}. Si el equipo se niega a arrancar desde
el USB con un mensaje como \enquote{Secure Boot violation} o \enquote{imagen no
autorizada}, hay que desactivarlo:

\begin{pasos}
  \item Apaga el equipo. Enciéndelo y pulsa varias veces la tecla de la
        configuración (suele ser \keys{F2}, \keys{Supr} o \keys{Esc}; ver el
        apéndice~\ref{cap:b-teclas-arranque}).
  \item Busca la pestaña \emph{Security} o \emph{Boot} y la opción
        \emph{Secure Boot}. Ponla en \emph{Disabled}.
  \item Guarda y sal con \keys{F10} y confirma con \emph{Yes}.
\end{pasos}

\begin{nota}
Desactivar el arranque seguro no daña el equipo ni borra nada. Si más
adelante vuelves a Windows, puedes activarlo otra vez.
\end{nota}

\section{Arranca el equipo desde el USB}\label{sec:arrancar-usb}

\begin{pasos}
  \item Con el equipo apagado, conecta la memoria USB. En una PC de
        escritorio, mejor en un puerto de atrás.
  \item Enciéndelo y, apenas aparezca el logo de la marca, pulsa
        repetidamente la tecla del \textbf{menú de arranque}. Las más comunes
        están en la tabla~\ref{tab:teclas-cortas}; la lista completa, en el
        apéndice~\ref{cap:b-teclas-arranque}.
  \item En la lista que aparece, elige el USB con las flechas (puede
        aparecer con su marca, como \emph{SanDisk} o \emph{Kingston}, o como
        \emph{USB HDD}) y pulsa \keys{Enter}. Si aparece dos veces, una con
        \emph{UEFI} delante, elige esa.
\end{pasos}

\begin{table}[htbp]\centering
\caption{Tecla del menú de arranque en las marcas más comunes}\label{tab:teclas-cortas}
\begin{tabular}{ll}
\toprule
Marca & Tecla habitual \\
\midrule
Acer, Dell, Lenovo, Toshiba & \keys{F12} \\
Asus & \keys{Esc} o \keys{F8} \\
HP y Compaq & \keys{Esc} y luego \keys{F9} \\
Sony Vaio & \keys{F11} o el botón \emph{Assist} \\
PC de escritorio & \keys{F8}, \keys{F11} o \keys{F12} \\
\bottomrule
\end{tabular}
\end{table}

Si en la lista no aparece el USB, entra a la configuración del equipo
(sección anterior), busca la sección \emph{Boot} y pon el USB primero en
el orden de arranque.

\subsection{El menú de Lázaro}
Si todo va bien, aparece un menú con dos opciones (escritas sin tildes,
porque ese menú no las muestra bien):
\begin{description}
  \item[Lazaro Linux] la normal. Se elige sola a los pocos segundos.
  \item[Lazaro Linux (modo grafico seguro)] úsala solo si con la primera la
    pantalla se queda en negro. Arranca con una imagen más básica que
    funciona con casi cualquier tarjeta de video.
\end{description}

\captura{menu-arranque}{El menú de arranque de Lázaro en un equipo con UEFI}

Si el USB no arranca, o se detiene con un error, revisa la
sección~\ref{sec:prob-no-arranca} de \emph{Solución de problemas}.

\fuentes
